\documentclass[11pt,a4paper]{article}
% Spolecny styl pro ucitelske pripravy MAT (Work/MAT)
% Verze 1.1, 2026-10-23
\NeedsTeXFormat{LaTeX2e}
\ProvidesPackage{priprava}[2026/10/23 v1.1 Sablona ucitelskych priprav MAT]

% Zakladni typografie a cestina
\RequirePackage[T1]{fontenc}
\RequirePackage[utf8]{inputenc}
\RequirePackage[czech]{babel}
\RequirePackage{lmodern}
\RequirePackage{microtype}
\RequirePackage{amsmath,amssymb}
\RequirePackage{xcolor}
\RequirePackage{enumitem}
\RequirePackage[most]{tcolorbox}
\RequirePackage{tabularx,booktabs,array}
\RequirePackage{fancyhdr}
\RequirePackage{needspace}
\RequirePackage{geometry}
\RequirePackage[hidelinks]{hyperref}

\geometry{a4paper,top=17mm,bottom=18mm,left=23mm,right=23mm}
\definecolor{notegray}{gray}{0.39}
\setlength{\parindent}{0pt}
\setlength{\parskip}{2.5pt plus 1pt minus 0.5pt}
\setlist[itemize]{topsep=2pt,itemsep=1pt,parsep=0pt,leftmargin=1.45em}
\setlist[enumerate]{topsep=2pt,itemsep=2pt,parsep=0pt,leftmargin=1.65em}
\setcounter{secnumdepth}{0}
\renewcommand{\arraystretch}{1.12}

% Zahlavi: vlevo rocnik | globalni tema, vpravo cislo: nazev hodiny
% Zapati zustava zcela prazdne.
\pagestyle{fancy}
\fancyhf{}
\renewcommand{\headrulewidth}{0.35pt}
\renewcommand{\footrulewidth}{0pt}
\setlength{\headheight}{14pt}
\newcommand{\nastavPripravu}[4]{%
  \fancyhead[L]{\small #1\enspace\textbar\enspace #2}%
  \fancyhead[R]{\small #3: #4}%
}

% Minimalisticke cernobile boxy. Kazda definice zvlast.
\tcbset{commonlesson/.style={enhanced,breakable=false,colback=white,colframe=black,
  colbacktitle=white,coltitle=black,boxrule=0.45pt,
  arc=0.8mm,outer arc=0.8mm,left=3mm,right=3mm,top=1.6mm,bottom=1.6mm,
  toptitle=1mm,bottomtitle=0.7mm,fonttitle=\bfseries,
  attach boxed title to top left={xshift=2.5mm,yshift=-1.15mm},
  boxed title style={colback=white,colframe=white,boxrule=0pt,left=1mm,right=1mm}}}
\newtcolorbox{zadaniBox}[1][]{commonlesson,title={Zadání},#1}
\newtcolorbox{prikladBox}[1][]{commonlesson,title={Příklad},#1}
\newtcolorbox{definiceBox}[1][]{commonlesson,title={Definice},#1}

% Spolecny box pro zadani a jeho reseni, bez mezery mezi castmi.
% Pouziti:
%   \begin{prikladReseniBox}
%     Zadani ulohy.
%     \tcblower
%     \begin{align*} ... \end{align*}
%   \end{prikladReseniBox}
% Spodni cast ma jemny sedy podklad a delici linku; lze v ni sazet
% vice radku matematiky. Obe casti maji jen jeden spolecny ramecek.
\tcbset{solutionlesson/.style={
  commonlesson,
  enhanced jigsaw,
  breakable,
  colbacklower=black!3,
  segmentation style={solid,draw=black!25,line width=0.35pt},
  fontlower=\small\color{notegray},
  before lower={\noindent\textbf{Řešení:}\enspace},
  leftlower=3mm,rightlower=3mm
}}
\newtcolorbox{prikladReseniBox}[1][]{solutionlesson,title={Příklad},#1}
\newtcolorbox{zadaniReseniBox}[1][]{solutionlesson,title={Zadání},#1}

% Pomucky pro autora pripravy: sedive poznamky; \reseni pro starsi pripravy
\newcommand{\poznamka}[1]{\par\begingroup\small\color{notegray}\itshape
  \textnormal{\textbf{Poznámka pro učitele.}}\enspace #1\par\endgroup}
\newcommand{\reseni}[1]{\par\begingroup\small\color{notegray}
  \textbf{Očekávané řešení:}\enspace #1\par\endgroup}
\newcommand{\blok}[1]{\needspace{5\baselineskip}\subsection*{#1}}
\newcommand{\vyrok}[1]{\uv{#1}}

\nastavPripravu{1. ročník}{Algebraické výrazy}{45}{Mocnění mnohočlenů}
\begin{document}

% Zdroje: vlastni ucebnice Algebraicke vyrazy, Umocnovani mnohočlenu;
% Realisticky.cz 030710--030712 Umocnovani mnohoclenu; 010802.
\poznamka{Jádro 45 minut: význam mocniny 5 min, druhá mocnina roznásobením 12 min, samostatná práce 12 min, třetí mocnina 10 min, shrnutí 6 min. Vzorce nezaváděj jako hotová pravidla; patří do hodiny 46.}
\blok{Co přesně znamená mocnina závorky?}
Mocnina znamená opakovaný součin \textbf{celého} výrazu, nikoli umocnění jednotlivých členů. Výpočty proto navazují na násobení mnohočlenů.
\begin{definiceBox}[title={Mocnina mnohočlenu}]
Pro mnohočlen $P$ a $n\in\mathbb N$ definujeme
\[P^n=\underbrace{P\cdot P\cdots P}_{n\text{ činitelů}},\qquad P^0=1.\]
\end{definiceBox}
\begin{prikladReseniBox}[title={Příklad 1: druhá mocnina dvojčlenu}]
Uprav $(x+3)^2$ bez použití vzorce.
\tcblower
\begin{align*}
(x+3)^2&=(x+3)(x+3)\\
&=x^2+3x+3x+9=x^2+6x+9.
\end{align*}
\end{prikladReseniBox}
\begin{prikladReseniBox}[title={Příklad 2: co udělá záporné znaménko?}]
Uprav $(2a-1)^2$ a $(1-2a)^2$. Porovnej výsledky.
\tcblower
$(2a-1)(2a-1)=4a^2-4a+1$ a $(1-2a)(1-2a)=1-4a+4a^2$. Výsledné mnohočleny jsou stejné: obě závorky jsou opačné, ale jejich druhé mocniny se rovnají.
\end{prikladReseniBox}
\poznamka{Jako typickou chybu předlož $(a+b)^2=a^2+b^2$ a ověř ji pro $a=b=1$: vlevo $4$, vpravo $2$. Nejde o platnou identitu. Vzorec s prostředním členem teprve odvodíme příště.}
\blok{Samostatná práce — druhá mocnina}
\begin{zadaniReseniBox}[title={S1: rozepiš jako součin a uprav}]
\begin{enumerate}[label=\alph*)]
\item $(x+2)^2$;
\item $(x-4)^2$;
\item $(2x+y)^2$;
\item $(3a-2b)^2$;
\item $(x^2+x)^2$.
\end{enumerate}
\tcblower
\begin{enumerate}[label=\alph*)]
\item $x^2+4x+4$;
\item $x^2-8x+16$;
\item $4x^2+4xy+y^2$;
\item $9a^2-12ab+4b^2$;
\item $x^4+2x^3+x^2$.
\end{enumerate}
\end{zadaniReseniBox}
\blok{Třetí mocnina: postupuj po etapách}
\begin{prikladReseniBox}[title={Příklad 3: třetí mocnina dvojčlenu}]
Uprav $(x+1)^3$ jako součin tří závorek.
\tcblower
\begin{align*}
(x+1)^3&=(x+1)(x+1)(x+1)\\
&=(x^2+2x+1)(x+1)\\
&=x^3+3x^2+3x+1.
\end{align*}
\end{prikladReseniBox}
\begin{prikladReseniBox}[title={Příklad 4: nejdřív umocni, pak dosazuj}]
Pro $P(x)=x-2$ urči $P(x)^3$ a hodnotu pro $x=3$.
\tcblower
$P(x)^3=(x-2)^3=x^3-6x^2+12x-8$; pro $x=3$ je hodnota $1$. Kontrola přímo: $(3-2)^3=1$.
\end{prikladReseniBox}
\blok{Volitelná rezerva}
\begin{zadaniReseniBox}[title={R1: mocnina trojčlenu}]
Uprav $(x+y+1)^2$ pomocí roznásobení.
\tcblower
$x^2+y^2+1+2xy+2x+2y$.
\end{zadaniReseniBox}
\blok{Domácí procvičení}
\begin{zadaniReseniBox}[title={D1}]
\begin{enumerate}[label=\alph*)]
\item $(a+5)^2$;
\item $(2x-3)^2$;
\item $(x-1)^3$;
\item $(x+y)^3$.
\end{enumerate}
\tcblower
\begin{enumerate}[label=\alph*)]
\item $a^2+10a+25$;
\item $4x^2-12x+9$;
\item $x^3-3x^2+3x-1$;
\item $x^3+3x^2y+3xy^2+y^3$.
\end{enumerate}
\end{zadaniReseniBox}
\end{document}
